% Concise block-preconditioner comparison with complete pMG--AMG coverage.
\providecommand{\ADRResultsPath}{docs/research/solver_studies/adr_scaling_2026_09_17}
\providecommand{\ADRResultsFigurePath}{run_outputs/solver_studies/adr_scaling_2026_09_17/figures}
\section{Block-preconditioned solvers for stationary advection--diffusion--reaction}
\label{sec:adr-polynomial-synthesis}

We compare block additive Schwarz (ASM) and block Jacobi (BJ), combined with
polynomial preconditioning of degree $d$, denoted ASM+PP($d$) and BJ+PP($d$),
against block algebraic multigrid (AMG) and face-block modal polynomial
multigrid with an AMGX coarse hierarchy ($p$MG--AMG). The family shorthand
$(\mathrm{ASM}\mid\mathrm{BJ})+\mathrm{PP}$ denotes either block
variant with a case-dependent polynomial degree; numerical labels give
the actual $d$. The finite-element degree $p_{\mathrm{FE}}$ is independent
of $d$. The comparison emphasizes convergence, iteration growth and solve
time as the mesh and degree increase.

ASM+PP, BJ+PP and $p$MG--AMG use the generalized minimal
residual method (GMRES). AMG denotes AMGX's classical block hierarchy,
with parallel modified independent-set coarsening, distance-two
interpolation, strength threshold $0.25$ and a direct coarse solve.
Its V-cycle has one pre/post sweep of BJ or multicolor block diagonal
incomplete lower--upper factorization (DILU), under flexible GMRES (FGMRES, F) or the
stabilized biconjugate-gradient method (BiCGSTAB, B).
$p$MG--AMG coarsens the symmetric part in modal coordinates on the same
mesh, with scalar AMG at degree zero and GMRES on the full operator.
Its standard policy (S) transfers directly to degree zero with degree-two
Chebyshev smoothing and one pre/post sweep; its robust policy (R) halves
the degree, with degree-four smoothing and two pre/post sweeps.

GPU measurements use 64-bit arithmetic on one RTX PRO 5000 Blackwell,
a zero initial guess, physical relative residual at most $10^{-10}$,
internal target $10^{-11}$ and a 1,000-iteration cap.
\emph{Fresh} time is the median setup-plus-first-solve time;
\emph{reused} time is the mean second-solve time with the same preconditioner.
Three independent setups follow one discarded trial (two setups for the
initial degree-four controls); timings cover only preconditioner setup and
linear solves, with validation outside the timer.
Each family retains its fastest passing configuration separately for
fresh and reused time. The associated $k$ is its mean outer iteration
count for that solve, rounded to an integer. Iterations of GMRES and
BiCGSTAB have different costs and are not interchangeable work units.
NC means no tested configuration in that family converged; a dash means
untested. Neither receives a ranked time.

\input{\ADRResultsPath/pardiso_lu/protocol.tex}

\subsection{Oscillatory cases favoring ASM+PP}
For the following classes, $\Omega$ is either $[-1,1]^2$ or the five-lobed
annulus $\{(r,\theta):0.58<r<1+0.35\cos(5\theta)\}$, with common fields
\begin{equation}
 \begin{aligned}
 u_\star&=\sin(6\pi x)\sin(5\pi y)+0.35\sin(11\pi x)\sin(9\pi y),\\
 \boldsymbol\beta_{\rm osc}&=\bigl(1+4\sin(7\pi x)\cos(7\pi y),\;
                   \tfrac12-4\cos(7\pi x)\sin(7\pi y)\bigr)^T.
 \end{aligned}
 \label{eq:adr-synth-common-osc}
\end{equation}
The square mesh ladder has 2,048, 8,192, 32,768 and 99,458 triangles;
the annular ladder has 2,043, 8,039, 33,174 and 99,984. The degree sweep
uses $p_{\mathrm{FE}}=1,2,3,4,6$ on each second mesh. Here $u_\star$ is the
manufactured exact solution, $f$ its matching source and $R(\theta)$ the
planar rotation through $\theta$. Polygonal boundaries receive the exact
Dirichlet data; the velocity need not be tangent to the walls.

\paragraph{Oscillatory transport dominance.}
\begin{equation}
 \begin{aligned}
 -10^{-5}\Delta u+\nabla\cdot(\boldsymbol\beta_{\rm osc}u)+10^{-2}u&=f,
 &u|_{\partial\Omega}&=u_\star.
 \end{aligned}
 \label{eq:adr-synth-osc-low}
\end{equation}
ASM+PP(24) is the fastest tested GPU method on both finest meshes, where both $p$MG--AMG
policies fail and the passing AMG configurations approach the iteration
cap. Its lower iteration counts support both the time advantage and the
larger convergence margin. Some lower-degree systems have no passing AMG
configuration, while ASM+PP(24) still converges.
The LU baseline is faster: on the finest annulus, $\PLU{16}$ takes
4.232\,s fresh and 0.689\,s reused, versus 13.773 and 12.369\,s for ASM+PP(24).

\paragraph{Weak anisotropy on the original annulus.}
\begin{equation}
 \begin{aligned}
 -\nabla\cdot(K_{\rm weak}\nabla u)+\nabla\cdot(\boldsymbol\beta_{\rm osc}u)+10^{-2}u&=f,
 &u|_{\partial\Omega}&=u_\star,\\
 K_{\rm weak}&=R(\pi/6)\operatorname{diag}(10^{-4},10^{-6})R(\pi/6)^T.
 \end{aligned}
 \label{eq:adr-synth-initial-weak}
\end{equation}
At 22,825 annular triangles and $p_{\mathrm{FE}}=6$, ASM+PP(24) takes
150 iterations and 0.759\,s reused, versus 230 iterations and 1.851\,s
for the fastest AMG configuration (FGMRES); both $p$MG--AMG policies fail.
Replacing $K_{\rm weak}$ by $10^{-5}I$ on this mesh leaves ASM+PP(24)
convergent in 300 iterations, while every tested AMG, stand-alone DILU
and $p$MG--AMG configuration fails. This control therefore retains an
ASM+PP(24) robustness advantage beyond its timing advantage.

\paragraph{Oscillatory directional anisotropy.}
\begin{equation}
 \begin{aligned}
 -\nabla\cdot(K_{\rm dir}\nabla u)+\nabla\cdot(\boldsymbol\beta_{\rm osc}u)+10^{-2}u&=f,
 &u|_{\partial\Omega}&=u_\star,\\
 K_{\rm dir}&=R(\pi/6)\operatorname{diag}(1,10^{-3})R(\pi/6)^T.
 \end{aligned}
 \label{eq:adr-synth-osc-directional}
\end{equation}
ASM+PP and $p$MG--AMG each converge on all 16 systems, whereas
all 32 tested AMG configurations fail. A stronger DILU-smoothed W-cycle
with three pre/post sweeps also fails at 85,120 unknowns: relative residual
$5.72$ after 1,000 FGMRES iterations. ASM+PP thus outperforms
these AMG configurations in robustness; $p$MG--AMG is faster at the
largest sizes.

\begin{table}[!htbp]
 \centering\small
 \resizebox{\linewidth}{!}{\input{\ADRResultsPath/synthesis/directional_tuned_table.tex}}
 \caption{Directional endpoints near one million unknowns at
 $p_{\mathrm{FE}}=6$. Apply is the mean duration of a separately profiled
 complete preconditioner application; setup is a median. Bold marks each
 geometry's minimum time in the last three columns. Here $k$ is the reused
 iteration count. LU cells give their thread count; CPU setup uses the
 fresh-selected count. LU has no preconditioner-application or Krylov count.}
 \label{tab:adr-synth-directional-tuned}
\end{table}

Iteration counts alone favor neither hierarchy uniformly. On the finest
square, ASM+PP(20) and $p$MG--AMG(S) have identical iteration and
application counts. On the annulus, ASM+PP(29) needs fewer of both, yet
$p$MG--AMG(S) is faster because its complete correction is much cheaper.
The robust policy reduces iterations but increases time on both geometries.
These measurements demonstrate a cost advantage of the tested
$p$MG--AMG correction, without establishing that ASM+PP is optimal.

\subsection{Crossovers under refinement and increasing degree}
The $(\mathrm{ASM}\mid\mathrm{BJ})+\mathrm{PP}$ vs AMG ranking
depends on problem size and on fresh versus reused timing. The following
cases separate iteration growth from the cost of each correction.
\paragraph{Oscillatory high diffusion.}
\begin{equation}
 \begin{aligned}
 -\Delta u+\nabla\cdot(\boldsymbol\beta_{\rm osc}u)+10^{-2}u&=f,
 &u|_{\partial\Omega}&=u_\star.
 \end{aligned}
 \label{eq:adr-synth-osc-high}
\end{equation}
Among GPU methods, ASM+PP(24) wins on coarse meshes, but its iteration count grows faster
than AMG's under refinement. On the annulus, the nearly constant
$p$MG--AMG(S) iteration count supports the lowest finest-mesh reused
time; its setup cost leaves AMG preferable fresh. Both geometry and
iteration growth therefore matter to the crossover.
At L4 on the annulus, AMG's 0.806\,s fresh and $p$MG--AMG's 0.133\,s
reused also beat $\PLU{24}$ (3.915 and 0.697\,s).

\begin{figure}[hbtp]
 \centering
 \includegraphics[width=\linewidth]{\ADRResultsFigurePath/synthesis/oscillatory_local_advantage.pdf}
 \caption{Oscillatory mesh scaling at $p_{\mathrm{FE}}=6$: reused time.
 Levels L1--L4 correspond to the mesh ladders above. Curve annotations
 give the polynomial degree, including the tuned directional endpoints;
 numbers beside the neutral LU curve give its MKL thread count.
 Crosses above the axes mark NC for the corresponding family.}
 \label{fig:adr-synth-oscillatory}
\end{figure}

\paragraph{Smooth anisotropic diffusion.}
\begin{equation}
 \begin{aligned}
 -\nabla\cdot\!\left(\operatorname{diag}(1,10^{-2})\nabla u\right)
 +(1,\tfrac12)^T\cdot\nabla u+u&=f,\\
 u|_{\partial[-1,1]^2}&=u_\star,\qquad
 u_\star=\sin(\pi x)\sin(\pi y).
 \end{aligned}
 \label{eq:adr-synth-smooth-anisotropic}
\end{equation}
ASM+PP(48) beats AMG throughout the fixed-mesh degree sweep. Under mesh
refinement its iteration count grows markedly, whereas AMG's stays nearly
constant and AMG becomes faster. $p$MG--AMG converges everywhere; at the
finest mesh, its robust policy reduces $k$ from 225 to 150 but increases
reused time from 1.03 to 2.52\,s. A stronger correction can cost more than
the iterations it saves.

\begin{figure}[hbtp]
 \centering
 \includegraphics[width=\linewidth]{\ADRResultsFigurePath/synthesis/smooth_anisotropic.pdf}
 \caption{Smooth anisotropic diffusion: reused time under mesh refinement
 and degree variation. Iteration counts and fresh times appear in
 Figures~\ref{fig:adr-synth-mesh-overview}--\ref{fig:adr-synth-degree-overview}.}
 \label{fig:adr-synth-smooth-anisotropic}
\end{figure}

\paragraph{Smooth transport dominance.}
\begin{equation}
 \begin{aligned}
 -10^{-3}\Delta u+(1,\tfrac12)^T\cdot\nabla u+u&=f,\\
 u|_{\partial[-1,1]^2}&=u_\star,\qquad
 u_\star=\sin(\pi x)\sin(\pi y).
 \end{aligned}
 \label{eq:adr-synth-smooth-transport}
\end{equation}
On the fixed mesh, BJ+PP(12) has lower fresh cost than ASM+PP(18), despite
requiring more iterations; ASM+PP(18) wins reused. Refinement favors AMG,
whose iteration growth is milder than BJ+PP(12)'s. At the million-unknown
endpoints, $p$MG--AMG remains convergent but needs far more iterations
than the polynomial methods. Stand-alone DILU is marked where it has
the lowest time.

\paragraph{Variable velocity.}
\begin{equation}
 \begin{aligned}
 -\Delta u+(1+x,-\tfrac12+y)^T\cdot\nabla u+\tfrac12u&=f,\\
 u|_{\partial[-1,1]^2}&=u_\star,\qquad u_\star=1+x^2+2y^2.
 \end{aligned}
 \label{eq:adr-synth-variable}
\end{equation}
BJ+PP(18) is attractive fresh on smaller systems; ASM+PP(18) usually
wins reused between these two methods. Their iteration counts rise with
refinement, while AMG and $p$MG--AMG remain much steadier. Consequently,
AMG minimizes finest-mesh fresh time and $p$MG--AMG minimizes reused
time, also at the $p_{\mathrm{FE}}=4$ million-unknown endpoint.

\paragraph{Smooth high diffusion.}
\begin{equation}
 \begin{aligned}
 -\Delta u+(1,\tfrac12)^T\cdot\nabla u+u&=f,\\
 u|_{\partial[-1,1]^2}&=u_\star,\qquad
 u_\star=\sin(\pi x)\sin(\pi y).
 \end{aligned}
 \label{eq:adr-synth-smooth-high}
\end{equation}
The same crossover holds: ASM+PP(18) and BJ+PP(12) are competitive on
small meshes, but BJ+PP(12) approaches the iteration cap under refinement.
AMG's iteration counts stay modest and it is fastest on the finest mesh.
$p$MG--AMG approaches AMG's reused time, but its setup cost prevents a
fresh-time advantage.

\begin{figure}[hbtp]
 \centering
 \includegraphics[width=\linewidth]{\ADRResultsFigurePath/synthesis/focused_h_crossover.pdf}
 \caption{Mesh refinement at $p_{\mathrm{FE}}=6$: fresh time above and
 reused time below. ASM+PP and BJ+PP remain separate.
 Stars mark an additional solver only when it is fastest overall.}
 \label{fig:adr-synth-focused-h}
\end{figure}

\begin{figure}[hbtp]
 \centering
 \includegraphics[width=\linewidth]{\ADRResultsFigurePath/synthesis/focused_p_crossover.pdf}
 \caption{Degree sweep on 8,192 triangles: fresh time above and reused time
 below, with the same configuration selection as the mesh sweep.}
 \label{fig:adr-synth-focused-p}
\end{figure}

\begin{table}[!htbp]
 \centering\scriptsize
 \setlength{\tabcolsep}{2.5pt}
 \resizebox{\linewidth}{!}{\input{\ADRResultsPath/synthesis/million_table.tex}}
 \caption{Near-million-unknown endpoints at $p_{\mathrm{FE}}=3,4$, on
 167,042 and 134,162 triangles, respectively. Each cell gives time in
 milliseconds, with configuration and $k$ below (LU gives threads, without $k$). Fresh and reused columns
 select configurations independently; bold marks the fastest time.}
 \label{tab:adr-synth-million}
\end{table}

\subsection{Practical interpretation}
The $(\mathrm{ASM}\mid\mathrm{BJ})+\mathrm{PP}$ family provides a
convergent GPU choice throughout the smooth and five-lobed cases: ASM+PP converges throughout,
including weak-diffusion oscillatory systems on which AMG and $p$MG--AMG
both fail. Directional anisotropy favors ASM+PP over the tested
AMG configurations in convergence, while $p$MG--AMG has cheaper large-system
corrections. On smooth problems, rising iteration counts limit
$(\mathrm{ASM}\mid\mathrm{BJ})+\mathrm{PP}$ under mesh refinement;
BJ+PP's smaller setup cost can nevertheless favor it fresh.
LU grounds these GPU comparisons: a gain over AMG does not by itself
establish an advantage over the CPU baseline. The nine-lobed stress
cases below further delimit iterative robustness.

\clearpage
\begin{figure}[hbtp]
 \centering
 \includegraphics[width=\linewidth]{\ADRResultsFigurePath/synthesis/solver_mesh_comparison.pdf}
 \caption{Mesh refinement at $p_{\mathrm{FE}}=6$: fresh and reused time
 with the corresponding iteration count. Each cell shows time in
 milliseconds followed by $(k)$, with the selected degree, Krylov method
 or policy below; LU instead gives $\PLU{n}$, with $n$ CPU threads.
 Color measures time relative to the fastest passing
 displayed family; its cell has a green border.}
 \label{fig:adr-synth-mesh-overview}
\end{figure}

\begin{figure}[hbtp]
 \centering
 \includegraphics[width=\linewidth]{\ADRResultsFigurePath/synthesis/solver_degree_comparison.pdf}
 \caption{Degree sweep on 8,192 square or 8,039 annular triangles, using
 the cell format of Figure~\ref{fig:adr-synth-mesh-overview}.
 Separate family outcomes preserve the ASM+PP result wherever
 AMG or $p$MG--AMG is NC.}
 \label{fig:adr-synth-degree-overview}
\end{figure}
\clearpage
